\clearpage\thispagestyle{empty}
\pdfbookmark[0]{Quellen zu Komet (21. September 1811)}{comet1811sept21-sources}
\begin{center}\textbf{Quellen zu Komet (21. September 1811)}\end{center}
\noindent Diese KI-Übertragung ins heutige Deutsch folgt dem Bericht in \href{https://gdz.sub.uni-goettingen.de/id/PPN236007467}{Gauß’ \textit{Werken}, Bd.~VI} (Göttingen, 1874), S.~337--338. Er erschien unter der Überschrift Göttingen in den \textit{Göttingischen gelehrten Anzeigen}, Stück 151 vom 21. September 1811, S.~\href{https://api.digitale-sammlungen.de/iiif/image/v2/bsb10538680_00473/full/1600,/0/default.jpg}{1497}, \href{https://api.digitale-sammlungen.de/iiif/image/v2/bsb10538680_00474/full/1600,/0/default.jpg}{1498} und \href{https://api.digitale-sammlungen.de/iiif/image/v2/bsb10538680_00475/full/1600,/0/default.jpg}{1499}. Das digitalisierte Exemplar stammt aus der Bayerischen Staatsbibliothek.
\par\smallskip\noindent Für den 22. Oktober gibt die Werkausgabe auf S.~338 die Stunde 0 an, der Erstdruck auf S.~1499 die Stunde 9. Der Bericht gibt die Abstände des Kometen von der Erde direkt an. KI-Übertragung und KI-Satz.
\clearpage
